\section{近期研究（2026 年 8--11 月）}
\paragraph{Swimm3R：先恢复水下几何，再联合拟合场景与介质}
2026.08~\cite{2608_00950}。普通 SfM 会把水体散射造成的对比度损失和颜色变化当成观测误差。Swimm3R 将 MASt3R-SfM 适配到水下图像，先估计恢复后的稠密点云、相机位姿与光学参数，再据此联合拟合水下 Beta 原语和显式介质场。散射感知的几何梯度与边缘感知的逆深度正则共同限制几何和水体效应的相互补偿。

在 FLSea 的几何评价里，相对作者列出的 MASt3R-SfM 基线，AbsRel 从 0.2247 降至 0.2132。对保留帧的渲染比较固定场景地图，但用各测试图像细化相机位姿 300 步，并对各方法沿用同一协议。部分位姿比较依赖重建参考轨迹，Beta 原语的渲染速度也低于标准三维高斯，宜分别看待位姿、保留帧画质和渲染成本。

\paragraph{QuerySplat：用独立查询分别回归几何与外观}
2026.08~\cite{2608_01186}。逐像素分配高斯会把原语数量绑定到输入视图与像素，而以场景查询统一回归所有属性又可能把空间结构与高频颜色混在一起。QuerySplat 使用双分支查询解码器，几何分支借助预训练的视觉几何模型提供坐标系、相机估计和三维特征，外观分支专门恢复辐射属性。输出仍是可直接泼溅的高斯，变化发生在高斯的分配和属性回归上。

论文在 DL3DV 的新视角合成比较中报告，相对该协议内最佳无位姿和需位姿基线，平均 PSNR 分别提高 2.30 dB 和 1.04 dB。这是针对论文所用视图和基线的比较。可扩展性仍需结合查询数、输入视图数和显存开销来判断。

\paragraph{StreamSplat：用三维缓存承接连续输入}
2026.08~\cite{2608_01659}。StreamSplat 用已标定的视图逐批构建三维高斯。在采用 DepthSplat 几何骨干与 ReSplat 式解码结构的基础上，方法将历史特征保存于有容量限制的场景坐标体素缓存。下一批输入到达时，缓存几何投影到当前视图以引导深度估计，邻近历史特征再影响高斯原语回归，因此每批输入后都能解码出可渲染场景。

原文在 DL3DV、RealEstate10K 和 ScanNet 上考察稀疏输入及长轨迹，对照 DepthSplat、AnySplat、ReSplat、Long-LRM 和使用官方权重的 OF³GS。长序列实验使用 256、512、1024 视图。方法针对静态场景并依赖位姿精度，不显式处理运动物体或光照改变。固定体素大小和缓存预算也在内存与几何细节之间形成取舍。

\paragraph{UniqueSplat：随查询视角调整前馈高斯}
2026.08~\cite{2608_02145}。固定高斯场从所有目标视角渲染同一组高斯。UniqueSplat 通过视角条件超网络，把目标视角相关的信息注入预测过程，在共享场景信息的基础上为每个查询生成适配的高斯。这给大角度视点变化提供了额外的调整自由度，也使表示的使用方式从一次预测后复用，变为由查询视角参与预测。

论文分别在 RealEstate10K、ACID 和 DTU 上评价新视角合成。作者在局限部分指出，高斯仍按输入视图像素对齐预测，未观测区域在渲染时仍可能出现伪影。这个限制要与已观测区域的画质改善分别评价。

\paragraph{InfiniSplat：从固定像素格改为沿表面采样高斯}
2026.08~\cite{2608_02437}。单图重建若在每个像素中心放置原语，较大基线下容易显露源相机采样的空隙。InfiniSplat 用单目深度或带测距输入的几何先验决定不规则支撑点，再向隐式解码器查询图像特征和高斯属性，并对初始高斯作有界更新。预测不需要逐场景拟合。

原文将 RGB 与 LiDAR 条件的模型分别同前馈基线比较，涉及 ETH3D、ScanNet++、Tanks and Temples 与 DL3DV。两种输入不能视作相同观测条件。论文展示了大面积反遮挡、错误深度先验和极端视角外推时的缺失或伪影，不能从一帧恢复完全不可见的结构。

\paragraph{FlexSplat：用变形查询汇聚未标定对象图像}
2026.08~\cite{2608_07937}。FlexSplat 对未标定的对象多视图联合训练几何变换器和高斯解码器，以预测深度建立空间锚定，再让多视图变形注意力汇聚证据。查询数限定高斯预算，与输入分辨率解耦，也不需要为每个输出原语建立点云对应。

作者在 GSO 将其与使用已知位姿的对象重建器比较，优势主要在感知画质与省去标定，PSNR 并非全面超过最强有位姿方法。模型训练与测试限于对象中心捕获，大场景背景和大的深度或位姿偏差仍会影响重建，平移扰动也有待完整评测。

\paragraph{LocusGS：让前馈查询在三维空间分工}
2026.08~\cite{2608_12825}。固定数量的全局查询能限制高斯输出预算，但一个查询生成的高斯可能散布在不相邻的区域。LocusGS 为每个输出查询维护三维中心和尺度，沿解码层更新。锚点既参与从多视图标定图像中选择相关证据，也约束该查询所解码高斯的局部范围，从而让不同查询负责互补区域。

作者在 RealEstate10K 的两视图设置，以相同查询数和高斯预算对照 TokenGS，在 PSNR、SSIM 和 LPIPS 上报告改善；DL3DV 另测从四视图训练到不同测试视图数的泛化。射线聚合仍依赖已标定相机，只有中心与标量半径的锚点也难以描述细长或高度各向异性的结构。

\paragraph{SplatGuide：从同一高斯重建提取三种生成条件}
2026.08~\cite{2608_16863}。缺少输入相机位姿时，二维视频生成先验难以独立维持跨视角几何。SplatGuide 用冻结的前馈重建骨干从静态场景参考图像估计位姿和高斯，再为新视图提供像素级渲染、特征 token 与可见性驱动的参考视图选择。最终渲染结果由扩散生成器细化，高斯在这里承担条件和几何接口。

论文区分域内和域外基准，并通过参考视图选择与条件消融评价各信号。三种条件来自同一重建骨干，因此纹理稀少、视角基线过宽或物体运动导致几何失准时，它们可能同时失效。目标视角按指定相机位姿查询，输入参考图无已知位姿这一条件不应与目标轨迹的定义混为一谈。

\paragraph{SeeU：用跨视图语义证据细化高斯}
2026.08~\cite{2608_22740}。SeeU 从带相机参数的稀疏图像估计初始高斯，再以跨视图熵感知的语义嵌入为条件，用 Transformer 在高斯空间回归残差。它改的是初始高斯的表示与细化阶段，输入图像依然提供场景观测约束。

作者在 RealEstate10K 和 ACID 的域内协议中用两张参考图渲染三个目标视角。RealEstate10K 训练后直接评估 DTU 时，SeeU 的 PSNR 为 16.31 dB，HiSplat 为 16.05 dB；在跨域 ACID 对比中，HiSplat 的 LPIPS 略优。两组结果适合分别讨论跨域几何与感知画质，不能据此推断到未知输入位姿。

\paragraph{利用逐视图高斯子集过滤瞬态干扰}
2026.08~\cite{2608_26951}。逐视图预测的前馈模型会为各输入视角生成相应的一组高斯。该工作利用这些集合的来源信息，将某个视图的高斯暂时排除，观察其余视图对候选重建内容的支持，再通过渲染检查删除候选是否能改善重建。模型权重保持不变，处理发生在推断后对候选高斯的筛选阶段。

论文在 RobustNeRF 与 NeRF On-the-go 上比较 4、8 和 16 个带干扰的输入视图，并检查干净场景中的画质。验证每个候选都需要渲染，因此计算开销随候选数和输入视图数增长；在多视角里一致出现的内容也不能仅靠单视图来源判为干扰。

\paragraph{ReconSplat：用可投影的高斯潜变量约束跨视图扩散}
2026.08~\cite{2608_28895}。稀疏宽基线输入中，可见结构与未观测区域需要不同的重建证据。ReconSplat 先从有位姿图像预测承载潜变量的高斯，再将这组三维潜变量投影到各目标视图，给多视角扩散提供一致的空间条件。模型同时支持 RGB 与深度外推，并从预测中恢复显式点云。

RealEstate10K 的内插与外推使用不同目标视图；论文报告外推中的多数指标优于所列基线，但 SSIM 例外。作者指出初始高斯的浮游伪影、缺失结构或不精确几何会传递到扩散条件。未观测区的生成结果可用于观察外推能力，其几何真实性还需要独立观测检验。

\paragraph{DualDiff3D：用结构与外观扩散先验迭代修正高斯}
2026.09~\cite{2609_01516}。宽基线高斯渲染的伪影既可能来自几何缺失，也可能来自外观细节不足。DualDiff3D 通过双分支扩散处理结构和外观，以结构感知注意力耦合两者，再执行渲染、修复和重新优化循环。置信度掩码用于选择修复结果可参与后续高斯优化的区域。

作者在 DL3DV 和 LLFF 分别评价逐图修复与最终高斯重建，并与常规 3DGS 和 DIFIX 系列比较。扩散器采用单步设计，但整个方法仍有额外算力开销，输入极度稀疏时也可能退化。因此应把单步修复和完整场景优化的耗时分开。

\paragraph{用三维高斯扰动构造视频修复训练对}
2026.09~\cite{2609_03657}。为学习能修复跨视图伪影的视频先验，通常要从多种欠观测条件反复拟合场景，再渲染训练视频。该工作直接扰动已拟合高斯的缩放、旋转和保留状态，生成带空间对应的退化视频，与同场景的干净渲染配对。免优化指构造退化场景的过程，视频先验仍需训练。

论文在多个稀疏观测设置比较二维退化和不同三维扰动，另用扩展模型评价深度与下游策略表现。三维扰动提供的是一种训练数据生成近似，对特定相机、退化分布和重建问题的覆盖仍应在目标数据上验证。

\paragraph{AVSplat：在多视图全局聚合之前引入辅助视图}
2026.09~\cite{2609_05925}。输入视图增多本应改善覆盖，但直接用全局注意力融合时，证据未必集中在相互支持的观测上。AVSplat 在全局聚合前让每个输入视图先从少量辅助视图获取场景上下文，再以占据情况调节体素融合权重，缓和致密输入下属性被平均得过于平滑的问题。

作者在 Mip-NeRF 360、VR-NeRF 和 DL3DV 比较稀疏与致密输入，报告相对所比较的无位姿前馈模型，增加视图时的结果更稳定。该结果说明输入规模利用率有改善，具体到显存和时延，还需按实际分辨率与输入量比较。

\paragraph{VoxelTTO：用体素聚合消除像素冗余，按相机先验做测试时校正}
2026.09~\cite{2609_21498}。像素对齐法会让不同视角为同一表面重复预测基元。VoxelTTO 将密集特征聚合到共享体素，用稀疏三维卷积和可分裂解码器回归高斯属性。论文还用随机实心体积渲染替换训练和推断中的常规高斯光栅化。

作者在 DTU、Replica 与 Tanks and Temples 上评估重建。测试时优化通过修改骨干网络中的轻量 LoRA 层，让预测相机与给定相机先验对齐。这一项依赖可用的相机监督，需要把纯前馈配置与带相机先验及测试时优化的配置分开评价。

\paragraph{GrapeSplat：在占用体素而非图像网格上生成高斯}
2026.09~\cite{2609_23182}。从每幅图固定位置解码高斯会随输入视图增加而产生重复原语。GrapeSplat 将无位姿图像的预测几何和特征提升到共同三维空间，在可压缩的体素坐标中汇总。稀疏解码器从占用体素生成初始高斯，使规模更多取决于占据空间，而非输入像素数。

论文在 NRGBD、7-Scenes、ScanNet++ 和 DL3DV 使用不同数量的上下文图像评测，并在附录扩展到 DTU 与 ETH3D。输入视图数变化和场景尺度变化是不同的泛化轴。补充材料的 DTU 示例仍有背景变体在相同区域失败，因此不能将体素汇总当成对所有上游几何错误的修正。

\paragraph{Spackle：把单图大视角补全交给高斯残差分支}
2026.09~\cite{2609_30941}。仅凭一张图渲染偏离很远的视角时，新露出的区域和原本可见的区域会争夺固定原语预算。Spackle 冻结预训练前馈模型，用渲染结果中的缺损掩码限定残差分支的训练区域，并把生成补全用作伪监督。推断时合并基础与残差高斯，无需再执行生成式补图。

论文在 Middlebury、Booster、WildRGBD、ETH3D、LLFF 和 Tanks and Temples 的大视点偏移协议下比较补全和原可见区域表现。新增区域的合理外观不代表恢复了真实遮挡面，补全质量受缺损掩码、伪监督与基础重建质量影响。

\paragraph{AESplat：将可见基本颜色与方向相关外观分开预测}
2026.09~\cite{2609_36693}。无位姿前馈高斯需要同时处理几何和外观。AESplat 直接用输入 RGB 构造球谐的零阶分量，只将高阶项交给带三维归纳偏置的预测模块。观测颜色提供基色，高阶分量负责补偿随方向变化的剩余项。

按原文对 NAS3R 的比较，RealEstate10K 与 ACID 的 PSNR 分别提高 0.803 与 0.847 dB。这个增益针对新视图合成。原文指出，完全不在输入视图中的内容难以可靠重建，用单次 RGB 近似方向无关颜色在强曝光差异或非朗伯反射下也有偏差。

\paragraph{StereoGaussians：从单对标定双目图像预测附近视角}
2026.09~\cite{2609_38592}。StereoGaussians 在冻结的立体网络中提取视差和多尺度特征，再用高斯预测头构建米制场景。第二层高斯与扩展画布给反遮挡和原视野外区域留出表达容量。SceneSplat-Stereo 用静态高斯教师渲染目标视图，提供成对训练监督。

论文在真实 StereoNVS 的八个留出场景与由已优化教师渲染的 ReplicaGS 分开评测；后者并非独立实拍目标。前层深度固定限制了对立体估计错误的修正，教师伪影仍会传入。原文评测集中在相近相机，不能推断大角度反遮挡效果。

\paragraph{3DROID：把机器人工作空间高斯锚定到米制坐标并单独测量可靠性}
2026.10~\cite{2610_01744}。视觉前馈模型生成的高斯可以合成图像，尺度与坐标却可能不适合机器人工作空间。3DROID 将精化相机外参与有位姿预测结合，并用覆盖感知的渲染深度、双目几何和机器人运动学表面一致性另外评价重建。相机外参提供对全局坐标的约束，度量几何精度仍需测量。

在 110 个场景的保留 ext2R 相机上，相对默认无位姿配置，YoNoSplat 的 PSNR 从 12.57 提升到 15.81 dB，ZipSplat 从 12.75 提升到 16.27 dB。该比较同时改变外参来源和位姿条件。数据目前以静态帧为主，几何评价限于可由立体匹配或机器人表面支持的位置，未纳入腕部相机。

\paragraph{LoCoSplat：把三维场景聚合简化为固定局部均值}
2026.10~\cite{2610_04351}。LoCoSplat 保留以深度和图像特征构造体素锚点的二维阶段，把学习式三维稀疏网络换成固定的局部特征平均，并通过逐点感知机与高斯预测头解码。高斯生成仍依赖二维阶段给出的深度和已知输入位姿，改动集中在锚点附近的三维上下文融合。

在 RealEstate10K 的 6、12、24 视图设置中，作者把方法与按相同协议重新训练的 VolSplat 对比。在 24 视图设置中，作者报告 30.57 dB，配对差为 3.3 dB。其他训练语料和分辨率的发布模型在表中单独标注，不能从该实验推断无位姿下的速度和画质。

\paragraph{DensiTok：在几何潜空间补齐未观测的中间相机}
2026.10~\cite{2610_07958}。前馈高斯骨干看到的真实视图很少时，DensiTok 在它的几何令牌空间增加输入。层级自适应 VAE 压缩各层几何令牌，相机条件扩散变换器一步预测缺失中间视图的潜码，再解码为骨干重建头能用的令牌。这避免先生成完整图像再重跑骨干编码器。

论文对 AnySplat、WorldMirror 和 Depth Anything 3 在 RealEstate10K 与 DL3DV 上比较画质、位姿、深度或运行时间。虚拟相机只在首末观测相机之间插值，论文也在这个区间取评测视图。由此得到的结果无法说明训练轨迹之外的视图外推表现。

\paragraph{DeltaSplat：把图像残差变成高斯属性的迭代更新}
2026.10~\cite{2610_09853}。DeltaSplat 用相机射线与当前渲染深度把二维残差定位到三维，学习逐步更新已有高斯的位置、不透明度与颜色。迭代是在前馈网络内运行的轻量细化模块，无需为每个场景反向优化，适合接在已有无位姿重建骨干之后。

作者在 DL3DV 的 140 场景测试集上比较 6、12 和 24 个输入视图。已知内参但无已知外参时，6 视图的 PSNR 为 26.64 dB，较所用前馈骨干提升 1.75 dB；论文另列全未标定设置。输入分辨率为 224 乘 224，高分辨率的计算成本和精度仍需单独评估。

\paragraph{CompactSplat：按局部纹理复杂度分配前馈高斯}
2026.10~\cite{2610_13177}。逐像素预测使平坦区域消耗大量相似高斯。CompactSplat 用纹理引导的区域划分配合分层金字塔采样，让复杂和简单区域采用不同密度，并允许推断时控制稀疏度。改进发生在原语空间分配，输出仍可由高斯渲染器消费。

论文在 RealEstate10K 的非像素对齐方法对照中报告，CompactSplat-L 使用约 72.6k 个高斯，LPIPS 为 0.117；TokenGS 用约 66k 个高斯且 PSNR 更高，LPIPS 为 0.147。画质随指标与预算变化。作者还指出划分依赖正交网格，主要消除单视角内冗余，跨视角的重复结构仍有压缩空间。

\paragraph{Holi-Spatial 360：用匹配约束稀疏全景重建并生成空间标注}
2026.10~\cite{2610_19502}。稀疏全景相距较远时，图像修复容易改变房间原有结构。Holi-Spatial 360 用 DA3 的匹配给高斯优化提供几何锚定，再用适配的图像与视频模型补充连续轨迹视图。重建和 RGB-D 验证还被用于生成三维实例、描述与空间问答，构成可用于空间推理的数据。

10 个 RealSee3D 场景的 LiDAR 对照中，各方法采用相同 DA3 初始化，2 厘米阈值的 F-score 从 DA3 的 0.2207 提高到 0.2617。该检验覆盖的是这些室内场景的几何，未观测区域的生成内容仍受对应关系可靠性和参考覆盖限制。

\paragraph{InfiniSplat v2：固定几何中心，用可学习属性筛掉低质量观测}
2026.10~\cite{2610_20341}。InfiniSplat v2 用冻结的 DA3 从图像和可选相机条件估计深度与位姿，再把点云转成高斯。前馈属性分支通过渲染监督学习每个点的贡献，后续用输入视图的渲染反馈调整属性。固定中心将几何定位与外观学习分开，学习的不透明度可抑制几何不一致的点。

论文分别评价有位姿和无位姿配置，报告六个数据集上平均 LPIPS 与 DISTS 的改善，并在 ETH3D 和 DTU 做消融。抑制错误高斯有助于画质，但固定中心无法修复严重的几何位置误差，像素对齐产生的重复原语和未观测空洞仍存在。

\paragraph{Anchor3DGS：以共享三维锚点替代逐像素生成}
2026.10~\cite{2610_21634}。Anchor3DGS 从未标定输入构造几何支架，以固定数量的三维锚点聚合可见多视图信息，在锚点之间交换局部结构上下文，再把锚点解码为局部高斯。锚点预算不再随像素或视图数直接增长，并通过熵引导将更多容量分配给复杂区域。

论文以 3 和 6 个输入视图在 VR-NeRF、MipNeRF360 与独立的 DL3DV 子集评价，在 MipNeRF360 与 DL3DV 上报告超过 2 dB 的 PSNR 增益。该数字对应论文所列的稀疏设置和基线；极少视图中的结构仍受支架质量、锚点分辨率和遮挡影响。

\paragraph{在无位姿前馈串联中校正坐标约定和尺度}
2026.10~\cite{2610_23039}。直接把相机估计网络的输出接到高斯预测器时，坐标定义或平移尺度不一致会引入误差。该工作保持 VGGT 和 TokenGS 冻结，在接口处转换坐标约定并做平移尺度校正。作者还单独比较从全部场景视图估计位姿和仅用稀疏上下文的情况。

在 140 个 DL3DV 场景上，全视图估姿配置报告 22.47 dB，四视图输入配置报告 22.44 dB，所复现的 COLMAP 条件为 22.98 dB。两种观测权限应分开列出，不能把全视图位姿结果直接归到四图重建协议。

\paragraph{FineSplat：用视图分组和三维分组保留小块细节}
2026.10~\cite{2610_25011}。细粒度图像块能保留细节，但直接对高分辨率多视图做全局注意力代价很高。FineSplat 先在少数视图组成的小组内处理 4 乘 4 像素图块，再把初始高斯渲染回参考视图，用残差反馈组织三维空间分组，让不同视图组的证据沿三维对应关系交互。

论文给出 DL3DV 的 6、12 和 24 视图结果及跨域比较。在无位姿评价中，论文冻结高斯并对每个目标相机位姿优化 200 步进行对齐；有真实位姿时不做额外对齐。这是测量画质的对齐协议，需与模型的前馈推断时间分开理解。

\paragraph{RLSplat：按预算直接生成高斯并限制解码成本}
2026.10~\cite{2610_27712}。在密集预测后剪枝只能降低最终高斯数，前端仍会为被删原语支付预测成本。RLSplat 的分配策略先确定各锚点分别生成多少高斯，再从指定总预算直接解码，以局部交叉注意力细化。预算因此同时作用于表示规模和解码成本。

作者在 ACID 按每视图 1k、2k 和 4k 个高斯与预算可控方法比较，在 8 视图、每视图 1k 个时相对 F4Splat 报告 1.22 dB 的 PSNR 优势。图像编码仍有固定开销，不能把解码成本随预算下降外推成整个系统同幅加速。实验还限于 256 乘 256 分辨率与一种跨域设置。

\paragraph{SG-Splat：在位姿估计有误时分阶段修正重建}
2026.10~\cite{2610_35361}。SG-Splat 将几何估计与高斯属性预测分开。模型先用光流引导细化位姿，再利用残差修正几何。双路径训练混合真实和预测位姿并使用不同的优化方式，使网络既能保留位姿监督，又能适应推断时的位姿误差。

论文在 DL3DV 和 RealEstate10K 比较稀疏与较密输入的新视角合成及相对位姿。该方法需要已知相机内参，高斯仍按像素对齐预测，输入视图重叠和较高分辨率都可能带来额外开销。

\paragraph{GeoGSplat：先以锚点预测空间结构，再在固定预算内解码高斯}
2026.10~\cite{2610_36760}。GeoGSplat 结合锚点与三平面特征作为有结构的场景表示，分层解码时先在锚点附近预测稠密三维点，再回归高斯。这样把固定容量的非像素对齐预测连接到显式几何，而无需让每张输入图都直接生成一个像素对齐的原语集合。

论文设置 4096 个锚点，每个锚点预测 16 个点，共 65536 个高斯，并在 DTU 与 NRGBD 评价合成和重建。这是固定的原语容量。比较不同视图数和复杂度的场景时，需要同时考虑预算能否覆盖稀少结构及训练和推断成本。

\paragraph{ThermalSplat：从无位姿稀疏热图预测高斯}
2026.11~\cite{2611_15435}。热图的纹理和辐射响应不同于 RGB。ThermalSplat 微调无位姿几何骨干，以软 Census 损失监督对亮度变化较稳健的局部结构，再用边界损失约束温度突变，从两张稀疏热图前馈得到可渲染高斯。

逐场景热重建基线的对照按各方法预设的保留视图计算指标，另一组泛化评估使用未见场景里所有可用评价视图。两种协议对应不同问题。论文明确假设静态场景，宽基线少图的未观测区域也仍有几何歧义。

\paragraph{PaveSplat：用残差旋转校正表面对齐的协方差}
2026.11~\cite{2611_26045}。沿预测表面对齐的解析协方差给出了合理的高斯形状初值，但局部估计还可能偏离真实切平面。PaveSplat 保留 SurfSplat 类的解析初值，将切面内的尺度与残差四元数用于局部校正。残差从单位旋转初始化，学习的变化集中在每个表面基元的协方差方向。

论文在 RealEstate10K 使用两张输入、256 乘 256 训练及标准测试划分，并在相同训练配置下比较不同渲染尺度。跨尺度表现与训练配方和表面深度质量有关，应与完全不同的数据域或原语预算分开比较。

\paragraph{FFGaussType：把已知字形轮廓转为可复用高斯资产}
2026.11~\cite{2611_26431}。FFGaussType 面向已知轮廓的挤出字形与 SVG 图标。方法从采样锚点预测二维特征平面，通过跨资产共享的解码器把平面特征转成带法向的神经高斯。相比逐资产拟合，学习的对象是从给定轮廓几何到高斯资产的转换。

论文按字体家族或图标来源拆分数据，报告全测试集法向误差 1.62 度，来自已准备锚点的转换耗时为 9.1 毫秒。该时间不含轮廓锚点的前处理。论文仅研究具有均一材质的挤出资产，着色不处理自阴影或与周边场景的光传输。
